% Generated draft: requires mathematical and source review.
\documentclass[11pt]{article}
\usepackage[margin=25mm]{geometry}
\usepackage{fontspec}
\setmainfont{Latin Modern Roman}
\usepackage{amsmath,amssymb,mathtools,unicode-math}
\setmathfont{Latin Modern Math}
\usepackage{enumitem,needspace,xurl,hyperref}
\setlength{\parindent}{0pt}
\setlength{\parskip}{5pt}
\setlength{\emergencystretch}{3em}
\newcommand{\N}{\mathbb{N}}
\newcommand{\Z}{\mathbb{Z}}
\newcommand{\Q}{\mathbb{Q}}
\newcommand{\R}{\mathbb{R}}
\newcommand{\C}{\mathbb{C}}
\newcommand{\F}{\mathbb{F}}
\newcommand{\A}{\mathbb{A}}
\newcommand{\PP}{\mathbb{P}}
\newcommand{\E}{\mathbb{E}}
\newcommand{\eps}{\varepsilon}
\newcommand{\dd}{\,\mathrm{d}}
\begin{document}
\section*{UnsolvedMath ID: 30000513}
\textbf{Source problem code:} catalogue entry
\subsection{Definitions and mathematical statement}
\textbf{Two generic-smoothness questions.} For a smooth curve \(X\) of general moduli and a general member \(V'\) of the Kuranishi family of \(V\), is every component of \(\operatorname{Hom}(X,V')\) generically smooth? Is every component of \(\operatorname{Hom}(X,G/P)\) generically smooth, with the same general-moduli condition on \(X\)? The two target classes are distinct questions and are retained separately.\par In the usual complex algebraic setting, \(\operatorname{Hom}(X,Y)\) is the parameter scheme of morphisms from the fixed curve \(X\) to the fixed target \(Y\), with its scheme structure, not merely its set of maps. The second target conventionally has \(G\) a complex reductive algebraic group and \(P\) parabolic, so \(G/P\) is a projective homogeneous variety. For an irreducible component \(Z\) of the underlying Hom space, the relevant meaning of generically smooth is that a dense open subset of \(Z\) consists of smooth points of the full Hom scheme. Passing first to a reduced component would discard the generic nilpotents at issue. In characteristic zero the generic smoothness of the reduced variety alone is not the assertion being asked.\par A Kuranishi family is a local deformation family of the intended compact complex target \(V\). The record does not give the hypotheses on \(V\), specify an irreducible component of its deformation base, or say whether the family and Hom object are algebraic schemes or complex analytic spaces. Smoothness and projectivity of the target must not be guessed if not stated in the missing context. Likewise, the genus of \(X\), the field and moduli space, general versus very general, and the order of generality quantifiers are omitted. Requiring a single general pair \((X,V')\) to work for every Hom component can differ from choosing a general pair separately for each fixed component or numerical class. No fixed degree or homology class is inserted into the universal component assertion.\par The supplied literature triage checked 2026-08-21 labels the record partially solved. It reports that Mukai and Nasu constructed a generically nonreduced component, giving a negative answer to the universal general-target formulation, while the corresponding homogeneous-target \(G/P\) question is explicitly left open. The example and its exact target, curve, and deformation hypotheses are not reproduced; these supplied claims require independent verification. In particular, the reported first counterexample is not automatically a counterexample for \(G/P\).
\subsection{Short English statement}
Ask for generic smoothness of every component of the full morphism space, first for a general deformation of a target and then for a homogeneous target. The supplied triage reports a negative general-target example but leaves the homogeneous case separate. Reducing the Hom scheme would erase the nonreducedness being tested.
\subsection{Sources}
\begin{itemize}\item Local UnsolvedMath registry, permanent ID 30000513, code catalogue entry; \emph{Generic Smoothness of Homomorphism Schemes}, Oberwolfach Reports Open Problems.\item \emph{Classical Algebraic Geometry (2007)}, Oberwolfach Reports, DOI 10.4171/owr/2006/27, pp. 1615--1662. \url{https://doi.org/10.4171/owr/2006/27}. Proposed year and year tag: 2006; no proposer named.\item Supplied literature label \emph{Mukai--Nasu, generalization of Mumford example}: \url{https://www.kurims.kyoto-u.ac.jp/~mukai/paper/owr06.pdf}.\item Supplied literature label \emph{original Oberwolfach report}: \url{https://doi.org/10.4171/owr/2006/27}.\item Literature triage checked 2026-08-21, status \texttt{partially\_solved}; record updated 2026-09-16. The \texttt{corrected\_verified} field reports source checking, not independently verified here.\end{itemize}
\end{document}
